% ── 4. System Architecture Diagram ───────────────────────────────────────────
\section{System Architecture Diagram}

% TODO: Export diagram from draw.io and place it below.

\subsection{AWS Service Descriptions and Responsibilities}

\subsubsection{Gateway Layer}

\textbf{AWS WAF + Application Load Balancer (ALB)} serves as the single entry point for all external traffic entering the system. The Web Application Firewall inspects inbound requests against managed rule sets to block common exploits such as SQL injection, cross-site scripting, and volumetric abuse before they reach any internal service. The Application Load Balancer sits behind the WAF and is responsible for TLS termination, protocol routing, and distributing requests across healthy target groups. MQTT traffic originating from device controllers is directed toward AWS IoT Core, while HTTP traffic from end users is split between the CloudFront-hosted frontend and the API Gateway, depending on the nature of the request. This layer ensures that no component in the internal infrastructure is directly reachable from the public internet.

\subsubsection{Connectivity Layer}

\textbf{AWS IoT Core} acts as the managed MQTT broker that maintains persistent, authenticated connections with all registered device controllers. It enforces per-device IoT policies that govern which topics a device may publish to or subscribe from, preventing cross-device interference or unauthorized command injection. IoT Core forwards device registration events to IoT Device Management, routes incoming video streams to Kinesis Video Streams, and publishes telemetry events to EventBridge for rule evaluation --- allowing device-sourced data to fan out to the appropriate downstream service without coupling the broker to individual consumer logic.

\subsubsection{Application Services}

\textbf{Amazon CloudFront + S3 (Frontend)} delivers the user-facing web application as a set of static assets cached at AWS edge locations globally. CloudFront reduces latency for geographically distributed users and offloads origin traffic from the backend, while S3 provides durable, low-cost object hosting for the compiled frontend bundle. Dynamic API calls made by the frontend are forwarded through CloudFront to API Gateway rather than being handled at the edge.

\textbf{Amazon Cognito} provides the user identity and authentication layer for the platform. It manages the consumer user pool --- handling account creation, credential storage, password policies, and multi-factor authentication --- and issues signed JWT tokens upon successful login. API Gateway validates these tokens on every inbound request before forwarding traffic to the backend, ensuring that unauthenticated sessions cannot reach any protected resource. Cognito also supports the invitation-based account linking flow required by the access delegation use case, where a secondary user accepts an invitation and either creates a new account or links an existing one to the inviting household.

\textbf{Amazon API Gateway} exposes the REST interface consumed by both the frontend application and any authorized third-party integrations. It handles request throttling, usage plan enforcement, and token validation via Amazon Cognito, ensuring that only authenticated sessions can invoke backend operations. API Gateway acts as the protocol boundary between HTTP clients and the ECS-hosted backend services, and it offloads cross-cutting concerns --- authentication, rate limiting, and CORS --- from the backend entirely.

\textbf{ECS Fargate (Backend)} runs the core application logic as containerized workloads without requiring management of the underlying compute infrastructure. Fargate tasks handle business operations including user account management, automation rule processing, device command dispatch, and orchestration of calls to downstream services. The backend is also responsible for registering rule-based schedules with EventBridge and for dispatching consumer-facing notifications through SNS in response to device events and household management operations such as access delegation invitations. Because Fargate scales task count independently of server provisioning, the backend can absorb traffic spikes driven by consumer growth without manual intervention.

\textbf{Amazon EventBridge} serves as the scheduling and event-routing backbone for the automation rules engine. When a Consumer defines a time-based automation rule, the backend registers a corresponding EventBridge schedule that fires at the specified cadence and invokes the backend rule evaluator as the target. EventBridge also receives device telemetry events forwarded from IoT Core and routes them to the appropriate backend handler for rule condition evaluation, decoupling the IoT message pipeline from direct backend invocations and enabling fan-out to multiple consumers of the same event stream.

\textbf{AWS IoT Device Management} provides the administrative plane for the controller fleet. It maintains the device registry, supports grouping of controllers by household or region, and serves as the channel through which over-the-air configuration and firmware update jobs are issued. Upon receiving a registration event from IoT Core, Device Management assigns the device an identity, associates it with the correct household record in Aurora, and then publishes the applicable device policy back through IoT Core so the controller begins operating within its authorized topic scope. This bidirectional relationship with IoT Core is what closes the registration handshake described in UC-04.

\textbf{Amazon Kinesis Video Streams} ingests, durably stores, and indexes live video data pushed from camera-equipped controllers via IoT Core. It handles the time-series nature of video data and exposes APIs for both real-time playback and historical retrieval --- the backend queries Kinesis Video Streams directly when a Consumer requests a live stream endpoint, receiving a time-limited URL that the frontend uses to initiate playback. Processed frames and media segments are ultimately archived to S3, where lifecycle policies can transition older footage to lower-cost storage tiers automatically.

\textbf{AWS Lambda (Media Collector)} provides event-driven processing of media artifacts as they arrive in S3. Lambda functions are triggered by S3 object creation events and are responsible for tasks such as thumbnail generation, metadata extraction, and transcoding. Critically, the Media Collector also writes media reference records --- including S3 object keys, timestamps, and associated device identifiers --- back to Aurora so that the historical event review use case can surface media clips alongside event log entries without requiring a direct scan of S3. Because Lambda scales to zero when idle, this component adds no standing cost during periods without media activity.

\textbf{AWS Backup} enforces a centralized, policy-driven backup strategy across both Aurora and S3. Backup plans define retention schedules, vault destinations, and cross-region copy rules, ensuring that recovery point objectives are met without relying on ad hoc scripting. In the event of data loss or corruption, AWS Backup provides a consistent restore path for both structured relational data and unstructured object data.

\subsubsection{Data Storage}

\textbf{Amazon Aurora (RDS)} stores all structured relational data for the system, including user accounts, household configurations, device registrations, automation rules, and audit logs. Aurora provides MySQL or PostgreSQL compatibility with a storage layer that automatically replicates across multiple Availability Zones and supports read replicas for query offloading. Its serverless scaling variant can also be considered to match compute capacity to actual query load during periods of low activity.

\textbf{Amazon S3} serves as the primary object store for all unstructured and semi-structured data, including video segments, media thumbnails, database backup archives, and frontend static assets. S3's lifecycle policies allow objects to transition automatically from standard storage to infrequent-access or Glacier tiers as they age, optimizing storage costs without requiring application-level logic to manage data tiering.

\subsubsection{Monitoring and Observability}

\textbf{Amazon CloudWatch} is the central observability service that collects metrics, logs, and traces from all internal components. ECS task metrics, IoT Core connection counts, Lambda invocation statistics, and Aurora performance data are all streamed to CloudWatch, providing a unified view of system health. CloudWatch Alarms evaluate metric thresholds and act as the trigger mechanism for automated responses and notifications.

\textbf{CloudWatch Dashboards} render the operational metrics collected by CloudWatch into visual panels accessible to engineering and operations teams. Dashboards can be scoped to individual services or composed into system-wide views, enabling teams to identify trends and anomalies during both normal operation and incident response.

\textbf{Amazon SNS} serves a dual notification role in the system. For consumer-facing use cases --- event alerts, automation triggers, and access delegation invitations --- the backend publishes directly to SNS topics scoped to individual households, which fan out to each Consumer's registered endpoints such as push notification tokens, SMS numbers, or email addresses. For operational use cases, CloudWatch Alarms publish to separate SNS topics that target the engineering on-call rotation via PagerDuty webhooks or similar integrations. Separating consumer and operational topic namespaces prevents alarm noise from reaching end users and ensures that consumer notification volume does not interfere with operational alerting reliability.

\textbf{Amazon QuickSight} connects to Aurora and CloudWatch data sources to produce business intelligence dashboards tracking key performance indicators such as active device counts, user engagement metrics, rule execution frequency, and system availability over time. QuickSight is intended for stakeholders and product teams rather than operations engineers, providing a business-level view of the platform's health and growth.
